% Paper 4, §10 — ROME and direct edits
% Thesis: Meng 2022 ROME bridges causal interpretability and weight-level
% model surgery. Causal tracing on GPT-2 XL localizes factual recall to
% mid-layer MLPs at the last subject token; rank-one MLP weight edits at
% exactly that locus install new (s, r, o) associations with 99.8% efficacy.
% Hase 2023 critique: locating != editing-target. The locus survives STR
% re-validation; the precise per-layer recommendation is corruption-sensitive.

\section{ROME and direct edits}
\label{sec:rome}

The activation-patching protocol of \S\ref{sec:patching} is a diagnostic:
its output is a heatmap of which $(\text{component}, \text{token},
\text{layer})$ cells are causal mediators for a behavior. ROME
\citep{meng2022-rome} is the first paper to take that diagnostic and
turn it into \emph{capability surgery} on the model's weights. The
argument has two halves. The first half \emph{localizes} factual recall
in GPT-2 XL with the three-run protocol of \S\ref{sec:patching-protocol}
and reports a sharp pair of peaks: an early site at mid-layer MLPs at
the last subject token, and a late site at the final-layer attention.
The second half \emph{intervenes}: treating the localized MLP as a
linear associative memory, ROME computes a closed-form rank-one weight
update that inserts a new $(s, r, o)$ triple --- ``the Eiffel Tower is
in Paris'' becomes ``the Eiffel Tower is in Rome'' --- with 99.8\%
efficacy and 88\% paraphrase robustness on the zsRE editing benchmark.
The methodological move is to use causal interpretability to identify
the locus, then weight-level surgery to validate the localization by
\emph{constructing} the predicted edit. The follow-up critique
\citep[Hase et al.\ 2023, ``Does Localization Inform Editing,''
forthcoming]{} \deepencite{hase2023-localization-editing §4,
citation-pending} complicates that move: the layers ROME's tracing
identifies are not always the layers where editing performs best.

\subsection{Causal tracing localizes factual recall}
\label{sec:rome-tracing}

\citet{meng2022-rome} apply the three-run protocol of
\S\ref{sec:patching-protocol} to factual prompts on GPT-2 XL,
sweeping the patched component $C$ over every $(\text{token}, \text{layer})$
cell of the residual-stream grid and computing the average indirect
effect $\mathrm{AIE}_C$ from
Eq.~\eqref{eq:aie}\footnote{KB excerpt:
\texttt{kb/excerpts/meng2022-rome.md\#sec-2-effects}.}. On 1000 factual
statements, two peaks dominate the
heatmap\footnote{KB excerpt:
\texttt{kb/excerpts/meng2022-rome.md\#sec-2-2}.}: a \emph{late site}
near the final-layer prediction (unsurprising; the model commits to
the answer there), and an \emph{early site} ``at the last token of
the subject name'' \citep[\S 2.2]{meng2022-rome}, which the authors
flag as ``a new discovery.'' Decomposing the early-site cells into
sublayer contributions sharpens the picture further: at the
last-subject-token early site, $\mathrm{AIE}_{\mathrm{MLP}} = 6.6\%$
versus $\mathrm{AIE}_{\mathrm{attn}} = 1.6\%$, with the AIE peak
across hidden states at $8.7\%$ at layer
$15$~\citep[\S 2.2]{meng2022-rome}. The hidden-state grid's strongest
mediators are MLPs in the middle of the stack, processing the last
subject token, well before any attention head has copied the answer
forward.

That localization motivates the paper's central conjecture, the
\emph{Localized Factual Association Hypothesis}~\citep[\S
2.3]{meng2022-rome}\footnote{KB excerpt:
\texttt{kb/excerpts/meng2022-rome.md\#sec-2-3}.}:

\begin{quote}
\itshape
each midlayer MLP module accepts inputs that encode a subject, then
produces outputs that recall memorized properties about that subject.
Middle layer MLP outputs accumulate information, then the summed
information is copied to the last token by attention at high layers.
\end{quote}

\noindent
The hypothesis localizes factual recall along three axes simultaneously
\citep[\S 2.3]{meng2022-rome}: (i) in the MLP sublayers, not the
attention sublayers; (ii) at specific middle layers; (iii) at the
processing of the subject's last token. It is the strongest-form
factual-localization claim the field has fielded, and the second half
of the paper validates it by intervention.

\subsection{The associative-memory view of MLPs}
\label{sec:rome-memory}

ROME's intervention rests on a specific reading of the MLP
sublayer~\citep[\S 3.1]{meng2022-rome}\footnote{KB excerpt:
\texttt{kb/excerpts/meng2022-rome.md\#sec-3-1}.}: the second linear
layer of an MLP, $W^{(\ell)}_{\mathrm{proj}}$, is a \emph{linear
associative memory} in the sense of \citet{kohonen1972-associative}
and Anderson 1972. Any linear map $W$ can act as a key--value store for
keys $K = [k_1 \mid k_2 \mid \ldots]$ and values $V = [v_1 \mid v_2 \mid
\ldots]$ by solving $W K \approx V$, with the Moore--Penrose
pseudoinverse minimizing squared error: $\hat{W} = V K^{+}$. In the MLP
context, $k$ is the gated activation produced by the first linear
layer (so $k$ is the MLP's encoded query for ``what subject is this''),
and $v = W^{(\ell)}_{\mathrm{proj}} k$ is the MLP output written into
the residual stream (so $v$ is the recalled information about the
subject).

Inserting a new association $(k_*, v_*)$ then has a closed-form
solution. Constrain the new map $\tilde{W}$ to satisfy
$\tilde{W} k_* = v_*$ exactly while leaving other key--value pairs
minimally perturbed. The closed-form rank-one update
\citep[\S 3.1, Eq.~2]{meng2022-rome} is
\begin{equation}
  \min_{\tilde{W}} \|\tilde{W} K - V\|
  \quad \text{such that} \quad
  \tilde{W} k_* = v_*,
  \quad
  \hat{W} = W + \Lambda (C^{-1} k_*)^\top,
  \label{eq:rome-edit}
\end{equation}
where $C = K K^\top$ is the uncentered covariance of keys (pre-cached
on a Wikipedia sample), and $\Lambda = (v_* - W k_*) / (C^{-1}
k_*)^{\top} k_*$ is a vector proportional to the residual error of the
new pair on the original memory~\citep[\S 3.1]{meng2022-rome}. The
update is rank one because $\Lambda (C^{-1} k_*)^{\top}$ is the outer
product of two vectors; it is closed-form because $C$ is a fixed
data-side statistic and $\Lambda$ is determined by the constraint.
Inserting any new fact reduces to (i) computing the key $k_*$ for the
target subject by running the model up to the chosen MLP layer,
(ii) computing the value $v_*$ that would produce the target object
$o$ at the unembed by gradient-descending against the LM loss with
the rest of the model frozen, and (iii) applying
Eq.~\eqref{eq:rome-edit}.

\begin{intuition}
The associative-memory view says: the MLP at the localized layer is
running the matrix multiply $v = W k$ as a literal lookup, with the
key $k$ being the model's internal representation of the subject and
the value $v$ being the residual-stream contribution that pushes the
unembed toward the correct object. ROME's edit replaces one row of the
lookup table without touching the rest. Returning to canonical math:
the constraint $\tilde{W} k_* = v_*$ in Eq.~\eqref{eq:rome-edit}
\emph{is} ``write a new row,'' and the rank-one structure
$\Lambda (C^{-1} k_*)^{\top}$ is the cheapest perturbation to $W$ that
satisfies that constraint without smashing the other rows --- the
$C^{-1}$ term reweights the update by which directions in key-space
are already crowded with existing associations, so the perturbation
prefers under-used directions.
\end{intuition}

\subsection{Empirical headline: editing zsRE}
\label{sec:rome-zsre}

The closed-form edit
of Eq.~\eqref{eq:rome-edit}, applied at GPT-2 XL's localized layer
$\ell^* = 15$ at the last subject token, achieves the editing
results in
Table~\ref{tab:rome-zsre}~\citep[\S 3.2]{meng2022-rome}\footnote{KB
excerpt: \texttt{kb/excerpts/meng2022-rome.md\#sec-3-2}.}.

\begin{table}[h]
\centering
\small
\begin{tabular}{@{}lccc@{}}
\toprule
Editor & Efficacy & Paraphrase & Specificity \\
\midrule
GPT-2 XL (no edit)  & 22.2 (\(\pm\)0.5) & 21.3 (\(\pm\)0.5) & 24.2 (\(\pm\)0.5) \\
Fine-tuning (FT)    & 99.6 (\(\pm\)0.1) & 82.1 (\(\pm\)0.6) & 23.2 (\(\pm\)0.5) \\
FT+L                & 92.3 (\(\pm\)0.4) & 47.2 (\(\pm\)0.7) & 23.4 (\(\pm\)0.5) \\
MEND                & 75.9 (\(\pm\)0.5) & 65.3 (\(\pm\)0.6) & 24.1 (\(\pm\)0.5) \\
\textbf{ROME}       & \textbf{99.8 (\(\pm\)0.0)} & \textbf{88.1 (\(\pm\)0.5)} & \textbf{24.2 (\(\pm\)0.5)} \\
\bottomrule
\end{tabular}
\caption{ROME on the zsRE editing benchmark, GPT-2 XL, replicated from
\citet[\S 3.2]{meng2022-rome}. \emph{Efficacy}: the edited model
produces the new $o$ on the exact training prompt. \emph{Paraphrase}:
the edited model produces $o$ on rephrased queries about the same
$(s, r)$. \emph{Specificity}: the edit does not contaminate predictions
on unrelated $(s', r)$ pairs.}
\label{tab:rome-zsre}
\end{table}

The methodological reading is what matters. ROME beats fine-tuning at
both efficacy (99.8 vs.\ 99.6) and paraphrase robustness (88.1 vs.\
82.1) while preserving specificity, with no gradient steps on the
broader model and no held-out training data --- only one rank-one
matrix update at the localized layer. That a closed-form,
parameter-frugal edit at the layer the patching diagnostic flagged
beats general-purpose fine-tuning is the substantive validation of
the Localized Factual Association Hypothesis: if the fact is stored
as a key--value pair at MLP layer $15$, you should be able to rewrite
it by rewriting that pair, and you can.

\subsection{Hase 2023: locating is not editing-target}
\label{sec:rome-hase-critique}

The headline that ``causal tracing localizes the layer to edit''
collapses in a more careful follow-up.
\deepencite{hase2023-localization-editing §1, citation-pending}
\emph{Does Localization Inform Editing} runs ROME-style edits at every
layer of GPT-J and GPT-2 XL, not only at the layer the tracing
diagnostic peaks at, and reports that the editing performance is
\emph{not monotonic} in the AIE peak: layers with substantially lower
$\mathrm{AIE}_{\mathrm{MLP}}$ frequently match or exceed the
editing performance at the tracing-localized layer, and the layer of
best edit is task-dependent in ways the tracing heatmap does not
predict. The methodological consequence is that \emph{causal tracing
identifies layers that are causal mediators of the behavior}, but
\emph{editing exploits a different property of the layer} --- something
closer to ``which layer's MLP, when perturbed, produces a residual-%
stream change that the rest of the network propagates faithfully to
the unembed for that fact.'' The two need not coincide, and on
Hase et al.'s benchmarks they often don't.

This is a falsification of the strong-form coupling claim --- causal
locus equals best edit-target --- but not of either half taken alone.
The localization claim survives: the tracing peak at GPT-2-XL layer
$15$ at the last subject token is robust, and we discuss its survival
under STR re-validation in \S\ref{sec:rome-corruption-sensitivity}.
The editing claim survives: rank-one edits at \emph{some} layer near
the localized band insert facts at 99.8\% efficacy. What does not
survive is the prescriptive coupling: ``compute the AIE heatmap, edit
at the peak.'' A more careful protocol selects the editing layer by
end-to-end editing performance on a held-out set, with the tracing
heatmap as a strong prior rather than a determinant.

\begin{contradiction}
\textbf{The ROME locus is robust; the layer-of-edit recommendation is
corruption-sensitive.} Two distinct claims travel together in the
ROME paper. The \emph{locus} claim --- ``mid-layer MLPs at the last
subject token are causal mediators of factual recall in GPT-2 XL''
--- survives STR re-validation: \citet[\S
3]{zhang2023-apatching}\footnote{KB excerpt:
\texttt{kb/excerpts/zhang2023-apatching.md\#sec-findings}.} report that
GN- and STR-derived heatmaps disagree on per-layer details but agree
on the broad early-site / late-site decomposition that the
hypothesis~\citep[\S 2.3]{meng2022-rome} rests on. The \emph{layer-of-edit}
claim --- ``edit at exactly the AIE peak, layer~15, for GPT-2 XL''
--- is more fragile. \citet[\S 1]{zhang2023-apatching} document that
GN-derived per-layer peaks shift under STR, and
\deepencite{hase2023-localization-editing §3, citation-pending}
document that editing performance does not track AIE peaks across
layers. The strongest defensible form of ROME's prescription, after
both critiques, is: \emph{causal tracing identifies a layer band where
factual recall is mediated; within that band, layer-of-edit should be
selected on a held-out editing benchmark, not by AIE peak alone}. The
locus is robust; the prescription is methodology-sensitive. Returning
to canonical math: the AIE statistic of Eq.~\eqref{eq:aie} is the
correct estimator for ``which components mediate the behavior''
\citep[\S 2.1]{zhang2023-apatching}, but it is not, by itself, an
estimator for ``which components, when perturbed by
Eq.~\eqref{eq:rome-edit}, produce edits that propagate
faithfully''---those are different objects and the field's experience
shows they decouple.
\end{contradiction}

\subsection{Frontier note: ROME at frontier-2026 architectures}
\label{sec:rome-frontier}

The ROME causal-tracing diagnostic and the rank-one edit of
Eq.~\eqref{eq:rome-edit} were developed on dense transformer MLPs in
GPT-2 / GPT-J. Whether the Localized Factual Association Hypothesis
extends to frontier-2026 architectures is partially open. Two
architectural shifts complicate the picture. First, MLA-compressed-KV
attention (DeepSeek-V2/V3) and grouped-query attention (Llama 3 family)
change the attention sublayer's computational structure, but the MLP
sublayers remain dense linear maps in those architectures, so the
associative-memory framing of \S\ref{sec:rome-memory} should carry
over by direct analogy. Second, fine-grained mixture-of-experts MLPs
(DeepSeek-V3 with 256 routed experts; Llama 4 MoE variants; Qwen 2.5
MoE) replace each MLP sublayer with a routing function over expert
modules, each of which is itself a dense linear map. The
associative-memory view applies per-expert, but the localization claim
splits: a fact is now stored across some subset of experts at the
localized layer, indexed by which experts the subject token routes to.
The closed-form rank-one edit of Eq.~\eqref{eq:rome-edit} would have
to be applied to the expert(s) the localized subject-token routes
through, with the routing distribution itself becoming a methodology
parameter. \deepencite{rome-on-frontier-2026 §full,
citation-pending}: a systematic causal-tracing study of factual recall
on Llama 3 / DeepSeek-V3 / Llama 4 with comparable benchmarks to
\citet{meng2022-rome} is, to our knowledge, not yet published; whether
the hypothesis survives the MoE / MLA architectural shifts is open
empirical.

\subsection{Cross-paper thread: editing as an alternative to fine-tuning}
\label{sec:rome-vs-finetune}

ROME's success makes weight-level editing a methodologically distinct
tool from fine-tuning for surgical knowledge updates. Paper~2's SFT
and preference-RL sections treat gradient-based fine-tuning (SFT,
RLHF, DPO) as the standard recipe for installing new behaviors and
knowledge; both update many or all parameters under a loss.
ROME-style editing is closed-form and updates a rank-one slice of one
matrix under a constraint. The trade-off is sharp: fine-tuning
generalizes to multi-fact updates and behavioral distributions; ROME
is cheap and surgical but, in its 2022 form, restricted to factual
$(s, r, o)$ triples and inherits the layer-of-edit fragility of
\S\ref{sec:rome-hase-critique}. LoRA/QLoRA adapter merging
\citep[Paper 2 §parameter-efficient-finetuning]{} is methodologically
adjacent --- a rank-$k$ update across many layers, gradient-trained
against a loss, where ROME is rank-one at one layer, closed-form
against a constraint. A unified axis (gradient-based vs.\ closed-form;
many-layer vs.\ one-layer; loss-driven vs.\ constraint-driven) is the
cross-paper synthesis; the substantive choice is task-dependent and
ROME-vs-fine-tuning at production scale is an open 2026 question.

\subsection{Forward link}
\label{sec:rome-forward}

ROME is the bridge between activation patching's diagnostic output
(\S\ref{sec:patching}) and the model-editing literature; it is also a
case study in the broader cross-method convergence picture
(\S\ref{sec:cross-method}). Section~\ref{sec:circuits} takes up the
generalization from a single localized component (mid-layer MLPs at
the last subject token) to a full circuit subgraph: where ROME's
factual-recall picture is essentially component-level (which sublayer,
at which token, at which layer), the circuit picture is
edge-level (which downstream component reads from which upstream
component, along which projection). Section~\ref{sec:cross-method}
returns to factual recall as a four-method convergence case ---
patching localizes the MLP; the tuned lens shows the predicted-object
probability rising across the localized layer band; SAE features on
those MLPs include subject-entity features; ROME's rank-one edit
succeeds at the localized layer --- treating it as the second canonical
\emph{positive} cross-method evidence after the IOI circuit.
